\section{Related Work}
\label{sec:related}

\paragraph{Vanilla speculative decoding.}
Leviathan \emph{et al.}~\cite{Leviathan2023} introduced the canonical
draft--verify scheme: a smaller \emph{draft} model proposes $k$ tokens
autoregressively and the larger \emph{target} verifies them in a single
parallel forward pass, with modified rejection sampling guaranteeing that the
output distribution is identical to that of the target. Chen
\emph{et al.}~\cite{Chen2023specsampling} concurrently proposed an equivalent
scheme. The expected per-cycle latency is
\begin{equation}
  L \;=\; \frac{T_{\text{draft}} + T_{\text{verify}}}{\tau},
  \qquad
  T_{\text{draft}} \;=\; k \cdot t_{\text{draft step}},
  \label{eq:specdec-latency}
\end{equation}
where $\tau \in [1, k{+}1]$ is the expected number of accepted tokens per
cycle. Wall-clock speedup $S = L_{\text{target}}/L$ therefore depends on the
draft/target cost ratio and the empirical token acceptance rate $\alpha$.
Stern \emph{et al.}~\cite{Stern2018blockwise} earlier explored blockwise
parallel decoding within a single model.

\paragraph{Reducing drafting cost.}
Medusa~\cite{Cai2024Medusa} eliminates the external draft by attaching
multiple prediction heads to the target model and verifying the resulting
candidate tree in one forward pass. The EAGLE
family~\cite{Li2024Eagle,Li2024Eagle2,Li2025Eagle3} keeps an external draft
but conditions it on the target's hidden features, reporting substantial
speedup under their respective settings. Despite these refinements, $T_{\text{draft}}$
remains linear in $k$, which forces EAGLE-3 to a single transformer layer
and caps achievable speedups at $\sim$$2$--$3\times$ on strong GPUs.

\paragraph{Diffusion-based drafters.}
Recent work attacks the linear-in-$k$ drafting bottleneck directly.
DiffuSpec~\cite{Li2025DiffuSpec} and SpecDiff-2~\cite{Sandler2025SpecDiff2}
explore diffusion-based parallel drafting to improve speculative-decoding
throughput. PARD~\cite{An2025PARD} instead focuses on low-cost parallel draft
model adaptation for accelerating LLM inference. DFlash~\cite{Chen2026dflash} closes this gap with a
lightweight (5-layer) \emph{block-diffusion} drafter conditioned on target
hidden features, reporting up to $5\times$ end-to-end speedup over
autoregressive baselines on H200/B200 hardware.

\paragraph{Precision and systems deployment.}
Deployment-oriented inference work shows that quantization and memory traffic
are central to practical throughput gains on constrained hardware.
Low-precision techniques such as LLM.int8()~\cite{Dettmers2022Int8} reduce
weight movement and can interact with speculative decoding in two opposing
ways: lower per-step cost for draft/verify passes, but possible shifts in
acceptance behavior when verifier and drafter precision settings differ.
This interaction motivates reporting speed and quality jointly rather than
throughput alone.

\paragraph{Distributed and edge relevance.}
In distributed and mobile computing settings, local inference can reduce
round-trip latency and cloud dependency for interactive workloads. Within this
context, speculative decoding is relevant not only to chatbot serving but also
to edge-side language interfaces in intelligent devices, where bounded latency
and predictable runtime behavior are as important as peak throughput.

\paragraph{Position of this work.}
This work focuses on a controlled consumer-GPU evaluation of vanilla
speculative decoding with same-family model pairing (Qwen2.5-3B target,
0.5B/1.5B drafts) under deterministic and stochastic regimes. Rather than
competing with data-center throughput claims, we characterise practical
operating points on a single RTX~4090 Laptop, including acceptance dynamics,
quality drift, and runtime-sensitive configuration choices. We then use these
results to motivate a future adaptive cascaded speculative decoding (ACSD)
extension aimed at reducing long-tail latency under the same hardware budget.
